\documentclass[10pt,a4paper]{amsart}
\usepackage[margin=19mm]{geometry}
\usepackage{amsmath,amssymb,mathtools}
\usepackage[scr=rsfs,cal=euler]{mathalfa}
\usepackage{xfrac}
\usepackage[unicode,hidelinks,bookmarks=false]{hyperref}
\newcommand{\ie}{i.e.\ }
\newcommand{\cf}{cf.\ }
\newcommand{\resp}{respectively\ }
\newcommand{\Subsection}{\subsection}
\providecommand{\nobreakdash}{\nobreak\hbox{-}}
\setlength{\emergencystretch}{2em}
\multlinegap0pt
\usepackage{enumerate}
\usepackage{amssymb}
\newtheorem{theorem}{Theorem}
\newtheorem{lemma}{Lemma}
\newcommand{\abs}[1]{\lvert#1\rvert}
\newcommand{\csch}{\operatorname{csch}}
\def\ge{\geqslant}
\title{Isoperimetric inequality for nearly spherical domains in the Bergman ball in $\mathbb{C}^n$}
\author{David Kalaj}
\date{Source: arXiv:2502.00891v3 (CC BY 4.0)}
\begin{document}
\maketitle

\noindent\emph{Source and licence.} Isoperimetric inequality for nearly spherical domains in the Bergman ball in $\mathbb{C}^n$. Version arXiv:2502.00891v3 (CC BY 4.0), distributed by arXiv under the Creative Commons Attribution 4.0 International licence, http://creativecommons.org/licenses/by/4.0/, as recorded in the arXiv licence metadata for this exact version and shown on the abstract page https://arxiv.org/abs/2502.00891v3. Full licence notice: \texttt{licenses/arxiv-2502.00891-CC-BY-4.0.txt}.\par\smallskip

\noindent\textit{Editorial scope.} Counted unit: the article's theorem mainth (source0.tex lines 403--430). The article's complete original proof of it is delivered with it (source0.tex lines 1008--1323, one \texttt{proof} environment of 316 lines, which the article places away from the statement). No mathematics is written, repaired or re-derived: the statement and the proof are the article's own lines, in the article's own notation, and no retained line is substituted or re-flowed.  Retained uncounted as the source-local context the proof needs: lines 324--326; lines 609--609; lines 611--611; lines 700--700; lines 783--787; lines 815--847; lines 1326--1344.  Nothing was omitted from any retained span, so the package carries no omission record.  No retained line contains a citation, so the wrapper carries no bibliography and no bibliography entry.  No retained line contains a figure environment, an image-inclusion command or drawing code, so no image is delivered and the package declares no asset dependency.  Editorial formatting only, in addition: an AMS \texttt{amsart} preamble that restates the article's own declarations and macros used by the retained lines, namely the environments \texttt{theorem}, \texttt{lemma} on separate counters, together with the standard packages those lines need.  The retained environments print in this document as Theorem 1, Lemma 1 in order of appearance, whereas the article prints the same items with the numbering of its own full document.  The complete native source of the article as distributed in the arXiv e-print for this exact version is bundled unchanged as \texttt{sources/172616/source0.tex}.
\begin{equation}\label{br}
P(\mathbb{B}_r) = \int_{\mathbb{S}_r}t^{2n-1}(1-t^2)^{-n} d\mathcal{H}(\omega),
\end{equation}

It follows that \begin{equation}\label{estU}|a_0|=\abs{\int_{\mathbb{S}}     u d\mathcal{H}(\omega)}\le c_0(n,r)\int_{\mathbb{S}} u^2 d\mathcal{H}(\omega)+ \varepsilon_0  C_1(n,r_0)\int_{\mathbb{S}} u^2 d\mathcal{H}(\omega).\end{equation}

\begin{equation}\label{rur}|a_0|\le \varepsilon_0 C_2(n,r_0)\|u\|_{W^{1,2}}.\end{equation}

\begin{equation}\label{estuom}\left|a_{1,i}\right|\le C_3(n,r_0)\varepsilon_0\|u\|_{W^{1,2}(\mathbb{S})}, i=1,\dots, 2n.\end{equation}

\begin{lemma}\label{lema22} Under previous notation, % Assume that $U(z) = \frac{2 \tanh^{-1}(|z|)}{1 + u\left(\frac{z}{|z|}\right)},$ $\mathbf{t}=\tanh[r/2(1+u(\omega))]$, $r=|z|$ and $\omega=z/|z|$.
if $0<r<r_0$ and $\|u\|_{W^{1,\infty}}<\varepsilon_0$ then
\[\begin{split}{\sqrt{1-\frac{|\left<\nabla U, z\right>|^2}{|\nabla U|^2}}}&\left[ 1 + \left( \frac{1 - \mathbf{t}^2}{2\mathbf{t}} \right)^2 r^2 |\nabla_\tau u|^2 \right]^{1/2}
\\&\ge 1+\frac{1}{2}  r^2 \csch^2 r |\nabla_\tau u|^2\\&+\frac{1}{8} r^2 \tanh^2(\frac{r}{2})(1-\tanh^2(\frac{r}{2}))(|\nabla_\tau u|^2-|\nabla_{\tau_{2n-1}} u|^2) \\&-\varepsilon_0 C_{8}(n,r_0)|\nabla_\tau u|^2.\end{split}\]
\end{lemma}

\begin{theorem}[Horizontal Poincar\'e inequality on $\mathbb S^{2n-1}$ and equality case]\label{prop23_general}
Assume that $n\ge 2$ and $U\in W^{1,2}(\mathbb{S}^{2n-1})$. Let
\[
U(z)=\sum_{k=0}^\infty\sum_{j=1}^{d_k} a_{k,j}\Phi_{k,j}(z),
\]
where $\{\Phi_{k,j}\}$ is an orthonormal basis of spherical harmonics on
$\mathbb S^{2n-1}$ satisfying
\[
-\Delta_{\mathbb S^{2n-1}}\Phi_{k,j}=k(k+2n-2)\Phi_{k,j}.
\]
Set
\[
TU(z):=\langle \nabla U(z),\imath z\rangle,
\qquad
U_k:=\sum_{j=1}^{d_k} a_{k,j}\Phi_{k,j}\in\mathcal H_k,
\]
where $\mathcal H_k$ denotes the space of spherical harmonics of degree $k$.
Then the following inequality holds:
\begin{equation}\label{neededbe_general}
\|\nabla U\|_{L^2(\mathbb S^{2n-1})}^2
-\|TU\|_{L^2(\mathbb S^{2n-1})}^2
\ge (2n-2)\sum_{k=0}^\infty\sum_{j=1}^{d_k} k\,a_{k,j}^2.
\end{equation}
Moreover, equality holds in \eqref{neededbe_general} if and only if for every $k\ge 1$
such that $U_k\not\equiv 0$ one has
\begin{equation}\label{eq:equality_condition_T}
TU_k = k\,U_k \quad \text{or}\quad TU_k = -k\,U_k
\qquad \text{on }\mathbb S^{2n-1}.
\end{equation}
Equivalently, each $U_k$ belongs to the $\pm k$--eigenspace of $T$ on $\mathcal H_k$.
In particular, equality holds precisely when $U$ is a sum of CR and/or anti-CR spherical
harmonics (no mixed bidegree components).
\end{theorem}

\begin{lemma}\label{lem:Anr0}
Let $n\ge 2$ and $r_0>0$. Define
\[
A_n(k,r):=
\frac12 r^2\csch^2(r)\,k(k+2n-2)
+\frac{2n-2}{8}\, r^2 \tanh^2\!\Big(\frac r2\Big)
\Big(1-\tanh^2\!\Big(\frac r2\Big)\Big)\,k,
\]
and
\[
d_0(n,r):=\frac{1}{2}r^{2}\left(n+(-1+n)\cosh r\right)\csch^{2}r,
\qquad
K_n(k,r):=A_n(k,r)-d_0(n,r).
\]
Then there exists a constant $A(n,r_0)>0$ such that for all $0<r\le r_0$ and all integers $k\ge 2$,
\begin{equation}\label{Kn_lower}
K_n(k,r)\ge A(n,r_0)\,\big(k(k+2n-2)+1\big).
\end{equation}
\end{lemma}

\begin{theorem}[Fuglede's theorem in the Bergman ball]\label{mainth}
For every $r_0>0$ there exists $\varepsilon_0\in\bigl(0,\tfrac12\bigr]$, depending only on $r_0$,
such that the following holds.

Let $E\subset \mathbb{B}_n$ be a set whose barycenter is at the origin and such that
\begin{equation}\label{eq:vol_constraint_thm}
\mu(E)=\mu(\mathbb{B}_r)
\end{equation}
for some $r\in(0,r_0]$. Assume that $\partial E$ is a radial graph over the unit sphere
$\mathbb{S}$, namely there exists a Lipschitz function $u:\mathbb{S}\to\mathbb{R}$ with
\begin{equation}\label{eq:graph_param_thm}
Z(\omega)=\omega\,\tanh\!\left(\frac r2\bigl(1+u(\omega)\bigr)\right),\qquad \omega\in\mathbb{S},
\end{equation}
such that $Z(\mathbb{S})=\partial E$, and
\begin{equation}\label{eq:u_small_thm}
\|u\|_{W^{1,\infty}(\mathbb{S})}\le \varepsilon_0.
\end{equation}
Then there exists a constant $c_1=c_1(r_0,n)>0$ such that
\begin{equation}\label{eq:fuglede_estimate}
\frac{P(E)-P(\mathbb{B}_r)}{P(\mathbb{B}_r)}
\;\ge\;
c_1\,\|u\|_{W^{1,2}(\mathbb{S})}^2.
\end{equation} In particular,
\[
P(E)\ge P(\mathbb B_r),
\]
with equality if and only if $E=\mathbb B_r$.
\end{theorem}

\begin{proof}[Proof of Theorem~\ref{mainth}]
Fix $n\ge 2$ and $r_0>0$. Let $E\subset\mathbb B_n$ be nearly spherical with barycenter at $0$
and assume $\mu(E)=\mu(\mathbb B_r)$ for some $r\in(0,r_0]$.
Let $u:\mathbb S^{2n-1}\to\mathbb R$ be Lipschitz with
$\|u\|_{W^{1,\infty}(\mathbb S^{2n-1})}\le \varepsilon_0$, where $\varepsilon_0$ will be chosen small.

\medskip

\noindent\textbf{Step 1: Set up the basic quantities.}
Set
\[
\mathbf t_\circ := \tanh\Big(\frac r2\Big),
\qquad
\mathbf t(\omega) := \tanh\Big(\frac{r(1+u(\omega))}{2}\Big).
\]
Define
\begin{equation}\label{G_general}
\begin{split}
G(r,\omega)
&:=\frac12 r^2\csch^2(r)\,|\nabla_\tau u|^2\\
&\quad+\frac18 r^2\tanh^2\!\Big(\frac r2\Big)\Big(1-\tanh^2\!\Big(\frac r2\Big)\Big)
\Big(|\nabla_\tau u|^2-|\nabla_{\tau_{2n-1}}u|^2\Big),
\end{split}
\end{equation}
where $\tau_{2n-1}$ denotes the Hopf direction on $\mathbb S^{2n-1}$ (i.e.\ the unit vector field
$T(\omega)= i\omega$).

By Lemma~\ref{lema22} and the expression \eqref{br} for $P(\mathbb B_r)$, we obtain
\begin{align}
P(E)-P(\mathbb B_r)
&\ge
\int_{\mathbb S^{2n-1}}\mathbf t^{2n-1}(1-\mathbf t^2)^{-n}
\Big(G(r,\omega)-\varepsilon_0 C_8(n,r_0)|\nabla_\tau u|^2\Big)\,d\mathcal H
\notag\\
&\quad+
\int_{\mathbb S^{2n-1}}
\Big(\mathbf t^{2n-1}(1-\mathbf t^2)^{-n}-\mathbf t_\circ^{2n-1}(1-\mathbf t_\circ^2)^{-n}\Big)\,d\mathcal H.
\label{Per_deficit_general}
\end{align}

\medskip

\noindent\textbf{Step 2: Taylor expansion of the Jacobian factor.}
Define the ratio
\[
J(u):=\frac{\mathbf t^{2n-1}(1-\mathbf t^2)^{-n}}{\mathbf t_\circ^{2n-1}(1-\mathbf t_\circ^2)^{-n}}.
\]
Then the second integral in \eqref{Per_deficit_general} is simply
\[
\mathbf t_\circ^{2n-1}(1-\mathbf t_\circ^2)^{-n}\int_{\mathbb S^{2n-1}}\big(J(u)-1\big)\,d\mathcal H.
\]

\smallskip

Since $\mathbf t(\omega)=\tanh\!\big(\frac r2(1+u(\omega))\big)$ and $|u|\le \varepsilon_0$,
we can expand $J(u)$ around $u=0$. More precisely, there exists a polynomial expansion
\begin{equation}\label{Jexp_general}
J(u)-1 = \alpha_1(r,n)\,u + \alpha_2(r,n)\,u^2 + R_{n,r}(w)\,u^3,
\end{equation}
where $w\in[-\frac12,\frac12]$ and the coefficients $\alpha_1,\alpha_2$ depend smoothly on $r,n$.
Moreover, the remainder satisfies
\begin{equation}\label{Rbound_general}
R_{n,r}(w)\,u^3 \ge -C_9(n,r_0)\,\varepsilon_0\,|u|^2
\qquad (0<r\le r_0).
\end{equation}
From \eqref{Jexp_general} and \eqref{Rbound_general} we obtain
\begin{equation}\label{Jlower_general}
J(u)\ge 1-C_7(n,r_0)\varepsilon_0.
\end{equation}

\medskip

\noindent\textbf{Step 3: Split the perimeter deficit.}
Divide \eqref{Per_deficit_general} by $\mathbf t_\circ^{2n-1}(1-\mathbf t_\circ^2)^{-n}$
and define
\[
\frac{P(E)-P(\mathbb B_r)}{\mathbf t_\circ^{2n-1}(1-\mathbf t_\circ^2)^{-n}}=:I_1+I_2,
\]
where
\begin{align*}
I_1
&:=\int_{\mathbb S^{2n-1}}J(u)\,
\Big(G(r,\omega)-\varepsilon_0 C_8(n,r_0)|\nabla_\tau u|^2\Big)\,d\mathcal H,
\\
I_2
&:=\int_{\mathbb S^{2n-1}}\big(J(u)-1\big)\,d\mathcal H.
\end{align*}

By \eqref{Jlower_general} and $J(u)\ge 0$, we obtain
\begin{equation}\label{I1lower_general}
I_1 \ge \int_{\mathbb S^{2n-1}}
\Big(G(r,\omega)-\varepsilon_0 C_8(n,r_0)|\nabla_\tau u|^2\Big)\,d\mathcal H.
\end{equation}

\medskip

\noindent\textbf{Step 4: Estimate $I_2$ using the volume constraint.}
The expansion \eqref{Jexp_general} yields
\[
I_2 = \alpha_1(r,n)\int_{\mathbb S^{2n-1}}u\,d\mathcal H
+\alpha_2(r,n)\int_{\mathbb S^{2n-1}}u^2\,d\mathcal H
+\int_{\mathbb S^{2n-1}}R_{n,r}(w)\,u^3\,d\mathcal H.
\]
The volume constraint $\mu(E)=\mu(\mathbb B_r)$ implies the estimate \eqref{estU}:
\[
\Big|\int_{\mathbb S^{2n-1}}u\,d\mathcal H\Big|
\le c_0(n,r)\int_{\mathbb S^{2n-1}}u^2\,d\mathcal H
+\varepsilon_0C_1(n,r_0)\int_{\mathbb S^{2n-1}}u^2\,d\mathcal H,
\]
where
\[
c_0(n,r)=\frac{r}{2\sinh r}\,(-1+n+n\cosh r).
\]
Combining this with the remainder estimate \eqref{Rbound_general} gives
\begin{equation}\label{I2lower_general}
I_2 \ge -d_0(n,r)\int_{\mathbb S^{2n-1}}u^2\,d\mathcal H
-C(n,r_0)\varepsilon_0\int_{\mathbb S^{2n-1}}u^2\,d\mathcal H,
\end{equation}
where $$d_0(n,r) = \alpha_2(n,r) - \alpha_1(n,r) c_0(n,r) = \frac{1}{2}r^{2}\left(n+(-1+n)\cosh r\right)\csch^{2}r.
$$
\medskip

\medskip
\noindent\textbf{Step 5: Coercivity via Lemma~\ref{prop23_general}.}
Recall that $\mathbb S^{2n-1}\subset\mathbb C^n$ carries the unit Hopf vector field
\[
T(\omega):= i\omega\in T_\omega\mathbb S^{2n-1}.
\]
We choose the orthonormal tangential frame $\{\tau_1,\dots,\tau_{2n-1}\}$ on $\mathbb S^{2n-1}$
so that
\[
\tau_{2n-1}(\omega)=T(\omega)=i\omega.
\]
In particular,
\[
\nabla_{\tau_{2n-1}}u = \langle \nabla u,\, i\omega\rangle.
\]

\smallskip
Let $u\in W^{1,2}(\mathbb S^{2n-1})$ and expand it in spherical harmonics:
\[
u(\omega)=\sum_{k=0}^\infty\sum_{j=1}^{d_k} a_{k,j}\,\Phi_{k,j}(\omega),
\qquad
-\Delta_{\mathbb S^{2n-1}}\Phi_{k,j}=k(k+2n-2)\,\Phi_{k,j},
\]
with $\{\Phi_{k,j}\}$ orthonormal in $L^2(\mathbb S^{2n-1})$.

\smallskip
Applying Lemma~\ref{prop23_general} with $U=u$, we obtain
\begin{equation}\label{coercivity_hopf}
\|\nabla u\|_{L^2(\mathbb S^{2n-1})}^2
-\|\langle\nabla u,i\omega\rangle\|_{L^2(\mathbb S^{2n-1})}^2
\ge (2n-2)\sum_{k=0}^\infty\sum_{j=1}^{d_k} k\,a_{k,j}^2.
\end{equation}
Since $\langle\nabla u,i\omega\rangle=\nabla_{\tau_{2n-1}}u$, this can be rewritten as
\begin{equation}\label{Hopf_difference}
\int_{\mathbb S^{2n-1}}
\Big(|\nabla_\tau u|^2-|\nabla_{\tau_{2n-1}}u|^2\Big)\,d\mathcal H
\ge (2n-2)\sum_{k=0}^\infty\sum_{j=1}^{d_k} k\,a_{k,j}^2.
\end{equation}

\smallskip
In particular, since the right-hand side vanishes for $k=0$ and equals $(2n-2)\sum_{j=1}^{d_1}a_{1,j}^2$
for $k=1$, we may isolate the high modes:
\begin{equation}\label{Hopf_high_modes}
\int_{\mathbb S^{2n-1}}
\Big(|\nabla_\tau u|^2-|\nabla_{\tau_{2n-1}}u|^2\Big)\,d\mathcal H
\ge (2n-2)\sum_{k\ge 2}\sum_{j=1}^{d_k} k\,a_{k,j}^2
\;+\;(2n-2)\sum_{j=1}^{d_1}a_{1,j}^2.
\end{equation}

\smallskip
We now substitute \eqref{Hopf_difference} into the definition of $G(r,\omega)$.
Recall that
\[
G(r,\omega)
=\frac12 r^2\csch^2(r)\,|\nabla_\tau u|^2
+\frac18 r^2 \tanh^2\!\Big(\frac r2\Big)\Big(1-\tanh^2\!\Big(\frac r2\Big)\Big)
\Big(|\nabla_\tau u|^2-|\nabla_{\tau_{2n-1}}u|^2\Big).
\]
Integrating $G$ over $\mathbb S^{2n-1}$ gives
\begin{align}
\int_{\mathbb S^{2n-1}}G(r,\omega)\,d\mathcal H
&=
\frac12 r^2\csch^2(r)\int_{\mathbb S^{2n-1}}|\nabla_\tau u|^2\,d\mathcal H
\notag\\
&\quad
+\frac18 r^2 \tanh^2\!\Big(\frac r2\Big)\Big(1-\tanh^2\!\Big(\frac r2\Big)\Big)
\int_{\mathbb S^{2n-1}}
\Big(|\nabla_\tau u|^2-|\nabla_{\tau_{2n-1}}u|^2\Big)\,d\mathcal H.
\label{intG_split}
\end{align}

\smallskip
Using the standard identity
\[
\int_{\mathbb S^{2n-1}}|\nabla_\tau u|^2\,d\mathcal H
=\sum_{k=0}^\infty\sum_{j=1}^{d_k} k(k+2n-2)\,a_{k,j}^2,
\]
together with \eqref{Hopf_difference}, we obtain the lower bound
\begin{equation}\label{G_spectral_lower}
\begin{split}
\int_{\mathbb S^{2n-1}} G(r,\omega)\, d\mathcal H
&\ge
\sum_{k=0}^\infty \sum_{j=1}^{d_k}
\left[
\frac12 r^2 \csch^2(r)\, k(k+2n-2)
\right.
\\
&\qquad\left.
+\frac18 r^2 \tanh^2\!\Big(\frac r2\Big)\Big(1-\tanh^2\!\Big(\frac r2\Big)\Big)
(2n-2)\, k
\right] a_{k,j}^2 .
\end{split}
\end{equation}


\smallskip
The important point is that for all $k\ge 2$ the coefficient in brackets is strictly positive,
and it grows quadratically in $k$, providing coercivity on the high-frequency part of $u$.
The low modes $k=0$ and $k=1$ are handled separately using the volume constraint and barycenter
conditions (see \eqref{rur} and the analogue of \eqref{estuom} in dimension $n$).
\smallskip
We now return to the term $I_1$. Recall from Step~3 that, using $J(u)\ge 1-C_7(n,r_0)\varepsilon_0$,
\begin{equation}\label{I1_basic_lower}
I_1 \ge \int_{\mathbb S^{2n-1}} G(r,\omega)\,d\mathcal H
- C_8(n,r_0)\varepsilon_0\int_{\mathbb S^{2n-1}}|\nabla_\tau u|^2\,d\mathcal H.
\end{equation}
Combining \eqref{I1_basic_lower} with \eqref{G_spectral_lower} and the identity
\[
\int_{\mathbb S^{2n-1}}|\nabla_\tau u|^2\,d\mathcal H
=\sum_{k=0}^\infty\sum_{j=1}^{d_k}k(k+2n-2)\,a_{k,j}^2,
\]
we obtain the following spectral lower bound for $I_1$:
we obtain the following spectral lower bound for $I_1$:
\begin{equation}\label{I1K_general}
I_1
\ge
\sum_{k=0}^\infty\sum_{j=1}^{d_k}A_n(k,r)a_{k,j}^2
- C_8(n,r_0)\varepsilon_0\|\nabla_\tau u\|_{L^2(\mathbb S^{2n-1})}^2.
\end{equation}
where
\[
A_n(k,r):=
\frac12 r^2 \csch^2(r)\,k(k+2n-2)
+\frac{2n-2}{8}\, r^2 \tanh^2\!\Big(\frac r2\Big)
\Big(1-\tanh^2\!\Big(\frac r2\Big)\Big)\,k .
\]


\medskip

\medskip
\noindent\textbf{Step 6: Lower bound the coercivity coefficient.}
From \eqref{I1K_general} and the estimate \eqref{I2lower_general}, and using the orthonormality of
$\{\Phi_{k,j}\}$ (so that $\int_{\mathbb S^{2n-1}}u^2\,d\mathcal H=\sum_{k,j}a_{k,j}^2$), we obtain
\[
I_1+I_2
\ge \sum_{k=0}^\infty\sum_{j=1}^{d_k} K_n(k,r)\,a_{k,j}^2
-C(n,r_0)\varepsilon_0\|u\|_{W^{1,2}}^2,
\]
where we define
\[\begin{split}
K_n(k,r)&:=
\frac12 r^2\csch^2(r)\,k(k+2n-2)
\\&+\frac{(2n-2)}{8}r^2\tanh^2\!\Big(\frac r2\Big)\Big(1-\tanh^2\!\Big(\frac r2\Big)\Big)\,k
-d_0(n,r).\end{split}
\]
Splitting off the low modes $k=0,1$ gives
\begin{equation}\label{I1I2_general}
I_1+I_2
\ge \sum_{k\ge 2}\sum_{j=1}^{d_k} K_n(k,r)\,a_{k,j}^2
-d_0(n,r)\Big(a_{0,0}^2+\sum_{j=1}^{d_1}a_{1,j}^2\Big)
-C(n,r_0)\varepsilon_0\|u\|_{W^{1,2}}^2.
\end{equation}

Then in Lemma~\ref{lem:Anr0} bellow, we prove  that there exists $A=A(n,r_0)>0$ such that for all $k\ge 2$,
\[
K_n(k,r)\ge A\big(k(k+2n-2)+1\big)
\qquad\text{for all }0<r\le r_0.
\]
(The proof follows the same monotonicity argument used for $n=2$ by considering
$H_n(k):=K_n(k,r)/(k(k+2n-2)+1)$ and checking it has a unique maximum, so the minimum
occurs at $k=2$ or $k\to\infty$.)

\medskip
\noindent\textbf{Step 7: Control of the low modes.}
The volume constraint gives \eqref{rur}, and the barycenter condition (the analogue of \eqref{estuom})
yields control of the $k=1$ coefficients. Hence there exists $C(n,r_0)>0$ such that
\[
a_{0,0}^2+\sum_{j=1}^{d_1}a_{1,j}^2
\le C(n,r_0)\varepsilon_0^2\,\|u\|_{W^{1,2}}^2.
\]

\medskip
\noindent\textbf{Step 8: Conclusion.}
Combining Steps 6 and 7 in \eqref{I1I2_general} and choosing $\varepsilon_0$ sufficiently small
(depending only on $n$ and $r_0$), we obtain
\[
I_1+I_2 \ge \frac{A}{2}\|u\|_{W^{1,2}(\mathbb S^{2n-1})}^2.
\]
Finally, by \eqref{br} the perimeter of the reference ball is
\[
P(\mathbb B_r)
=\omega_{2n-1}\,\mathbf t_\circ^{2n-1}(1-\mathbf t_\circ^2)^{-n},
\]
where $\omega_{2n-1}=|\mathbb S^{2n-1}|$. Therefore,
\[
D(E)=\frac{P(E)-P(\mathbb B_r)}{P(\mathbb B_r)}
=\frac{I_1+I_2}{\omega_{2n-1}}
\ge c_1(n,r_0)\|u\|_{W^{1,2}(\mathbb S^{2n-1})}^2,
\]
for some $c_1(n,r_0)>0$ depending only on $n$ and $r_0$.
This completes the proof.

\end{proof}

\end{document}
